% GENERATED FILE. Edit canonical lessons, then run npm run generate:books.
% canonical-source-hash: fnv1a64:d43b16a2089ae01c
% canonical-lessons: KA-C142-sundara, KA-C142-ala, KA-C142-agala, KA-C142-kiridu, KA-C142-odde

\chapter{Beautiful, Deep, Wide, Narrow, Wet}
\label{ch:ka-beautiful-deep-wide-narrow-wet}
\hlchaptermodality{142}

\begin{chapteropening}
\textbf{By the end of this chapter:} \emph{I can say beautiful, deep, wide, narrow and wet.}

Complete the last lesson of chapter 142: i can say beautiful, deep, wide, narrow and wet.
\end{chapteropening}

\section[\texorpdfstring{\kn{ಸುಂದರ} (sundara)}{ಸುಂದರ (sundara)}]{\kn{ಸುಂದರ} (sundara) — beautiful}

\label{lesson:KA-C142-sundara}



\begin{quote}
Before the new one: say the Kannada for dirty, then the Kannada for speed; fast.
\end{quote}

\subsection*{You'll want to know: \kn{ಸುಂದರ}}
\textbf{\kn{ಸುಂದರ}} — \emph{sundara} — \textquotedblleft{}beautiful\textquotedblright{}.

\begin{grammarlens}[title={What you've built}]
One more word for this chapter.
\end{grammarlens}

\subsection*{Your turn}
\begin{itemize}
\raggedright
  \item \emph{Say it:} \emph{sundara}
  \item \emph{Say it:} \emph{sundara}, once more
  \item \emph{Say it:} say \emph{sundara}
  \item \emph{Recall:} say the Kannada for dirty, then the Kannada for speed; fast, then say \emph{sundara} again
\end{itemize}

\begin{tcolorbox}[breakable,colback=teal!4,colframe=teal!35!black,title={Before you move on}]
What does \kn{ಸುಂದರ} mean? (\textquotedblleft{}Beautiful\textquotedblright{}.) Say it once more.
\end{tcolorbox}
\section[\texorpdfstring{\kn{ಆಳ} (āḷa)}{ಆಳ (āḷa)}]{\kn{ಆಳ} (āḷa) — deep; depth}

\label{lesson:KA-C142-ala}



\begin{quote}
Before the new one: say the Kannada for speed; fast, then the Kannada for beautiful.
\end{quote}

\subsection*{You'll want to know: \kn{ಆಳ}}
\textbf{\kn{ಆಳ}} — \emph{āḷa} — \textquotedblleft{}deep; depth\textquotedblright{}.

\begin{grammarlens}[title={What you've built}]
One more word for this chapter.
\end{grammarlens}

\subsection*{Your turn}
\begin{itemize}
\raggedright
  \item \emph{Say it:} \emph{āḷa}
  \item \emph{Say it:} \emph{āḷa}, once more
  \item \emph{Say it:} say \emph{āḷa}
  \item \emph{Recall:} say the Kannada for speed; fast, then the Kannada for beautiful, then say \emph{āḷa} again
\end{itemize}

\begin{tcolorbox}[breakable,colback=teal!4,colframe=teal!35!black,title={Before you move on}]
What does \kn{ಆಳ} mean? (\textquotedblleft{}Deep; depth\textquotedblright{}.) Say it once more.
\end{tcolorbox}
\section[\texorpdfstring{\kn{ಅಗಲ} (agala)}{ಅಗಲ (agala)}]{\kn{ಅಗಲ} (agala) — wide}

\label{lesson:KA-C142-agala}



\begin{quote}
Before the new one: say the Kannada for beautiful, then the Kannada for deep; depth.
\end{quote}

\subsection*{You'll want to know: \kn{ಅಗಲ}}
\textbf{\kn{ಅಗಲ}} — \emph{agala} — \textquotedblleft{}wide\textquotedblright{}.

\begin{grammarlens}[title={What you've built}]
One more word for this chapter.
\end{grammarlens}

\subsection*{Your turn}
\begin{itemize}
\raggedright
  \item \emph{Say it:} \emph{agala}
  \item \emph{Say it:} \emph{agala}, once more
  \item \emph{Say it:} say \emph{agala}
  \item \emph{Recall:} say the Kannada for beautiful, then the Kannada for deep; depth, then say \emph{agala} again
\end{itemize}

\begin{tcolorbox}[breakable,colback=teal!4,colframe=teal!35!black,title={Before you move on}]
What does \kn{ಅಗಲ} mean? (\textquotedblleft{}Wide\textquotedblright{}.) Say it once more.
\end{tcolorbox}
\section[\texorpdfstring{\kn{ಕಿರಿದು} (kiridu)}{ಕಿರಿದು (kiridu)}]{\kn{ಕಿರಿದು} (kiridu) — narrow}

\label{lesson:KA-C142-kiridu}



\begin{quote}
Before the new one: say the Kannada for deep; depth, then the Kannada for wide.
\end{quote}

\subsection*{You'll want to know: \kn{ಕಿರಿದು}}
\textbf{\kn{ಕಿರಿದು}} — \emph{kiridu} — \textquotedblleft{}narrow\textquotedblright{}.

\begin{grammarlens}[title={What you've built}]
One more word for this chapter.
\end{grammarlens}

\subsection*{Your turn}
\begin{itemize}
\raggedright
  \item \emph{Say it:} \emph{kiridu}
  \item \emph{Say it:} \emph{kiridu}, once more
  \item \emph{Say it:} say \emph{kiridu}
  \item \emph{Recall:} say the Kannada for deep; depth, then the Kannada for wide, then say \emph{kiridu} again
\end{itemize}

\begin{tcolorbox}[breakable,colback=teal!4,colframe=teal!35!black,title={Before you move on}]
What does \kn{ಕಿರಿದು} mean? (\textquotedblleft{}Narrow\textquotedblright{}.) Say it once more.
\end{tcolorbox}
\section[\texorpdfstring{\kn{ಒದ್ದೆ} (odde)}{ಒದ್ದೆ (odde)}]{\kn{ಒದ್ದೆ} (odde) — wet}

\label{lesson:KA-C142-odde}



\begin{quote}
Before the new one: say the Kannada for wide, then the Kannada for narrow.
\end{quote}

\subsection*{You'll want to know: \kn{ಒದ್ದೆ}}
\textbf{\kn{ಒದ್ದೆ}} — \emph{odde} — \textquotedblleft{}wet\textquotedblright{}.

\begin{grammarlens}[title={What you've built}]
One more word for this chapter.
\end{grammarlens}

\subsection*{Your turn}
\begin{itemize}
\raggedright
  \item \emph{Say it:} \emph{odde}
  \item \emph{Say it:} \emph{odde}, once more
  \item \emph{Say it:} say \emph{odde}
  \item \emph{Recall:} say the Kannada for wide, then the Kannada for narrow, then say \emph{odde} again
\end{itemize}

\begin{tcolorbox}[breakable,colback=teal!4,colframe=teal!35!black,title={Before you move on}]
What does \kn{ಒದ್ದೆ} mean? (\textquotedblleft{}Wet\textquotedblright{}.) Say it once more.
\end{tcolorbox}
